\documentclass[a4paper, 14pt]{extarticle}

\usepackage{tikz}
\usepackage{aeguill}

% Поля
%--------------------------------------
\usepackage{geometry}
\geometry{a4paper,tmargin=2cm,bmargin=2cm,lmargin=3cm,rmargin=1cm}
%--------------------------------------


%Russian-specific packages
%--------------------------------------
\usepackage[T2A]{fontenc}
\usepackage[utf8]{inputenc} 
\usepackage[english, main=russian]{babel}
%--------------------------------------

\usepackage{textcomp}
% § for list
\usepackage{easylist}

% Красная строка
%--------------------------------------
\usepackage{indentfirst}               
%--------------------------------------             


%Graphics
%--------------------------------------
\usepackage{graphicx}
\graphicspath{ {./images/} }
\usepackage{wrapfig}
%--------------------------------------

% Полуторный интервал
%--------------------------------------
\linespread{1.3}                    
%--------------------------------------

%Выравнивание и переносы
% \usepackage{tabularx}
%--------------------------------------
% Избавляемся от переполнений
\sloppy
% Запрещаем разрыв страницы после первой строки абзаца
\clubpenalty=10000
% Запрещаем разрыв страницы после последней строки абзаца
\widowpenalty=10000
%--------------------------------------

%Списки
\usepackage{enumitem}

%Подписи
\usepackage{caption} 

%Гиперссылки
\usepackage{hyperref}

\hypersetup {
    unicode=true,
    colorlinks=true,
    linkcolor=blue,
    filecolor=magenta,      
    urlcolor=cyan,
}

%Рисунки
%--------------------------------------
\DeclareCaptionLabelSeparator*{emdash}{~--- }
\captionsetup[figure]{labelsep=emdash,font=onehalfspacing,position=bottom}
%--------------------------------------

\usepackage{tempora}

%Листинги
%--------------------------------------
\usepackage{listings}

\lstset{
  basicstyle=\ttfamily\footnotesize, 
  %basicstyle=\footnotesize\AnkaCoder,        % the size of the fonts that are used for the code
  breakatwhitespace=false,         % sets if automatic breaks shoulbd only happen at whitespace
  breaklines=true,                 % sets automatic line breaking
  captionpos=t,                    % sets the caption-position to bottom
  inputencoding=utf8,
  frame=single,                    % adds a frame around the code
  keepspaces=true,                 % keeps spaces in text, useful for keeping indentation of code (possibly needs columns=flexible)
  keywordstyle=\bf,       % keyword style
  numbers=left,                    % where to put the line-numbers; possible values are (none, left, right)
  numbersep=5pt,                   % how far the line-numbers are from the code
  xleftmargin=25pt,
  xrightmargin=25pt,
  showspaces=false,                % show spaces everywhere adding particular underscores; it overrides 'showstringspaces'
  showstringspaces=false,          % underline spaces within strings only
  showtabs=false,                  % show tabs within strings adding particular underscores
  stepnumber=1,                    % the step between two line-numbers. If it's 1, each line will be numbered
  tabsize=2,                       % sets default tabsize to 8 spaces
  title=\lstname                   % show the filename of files included with \lstinputlisting; also try caption instead of title
  % Для отображения русского языка
  extendedchars=true,
  literate={Ö}{ {\"O} }1
  {Ä}{ {\"A} }1
  {Ü}{ {\"U} }1
  {ß}{ {\ss} }1
  {ü}{ {\"u} }1
  {ä}{ {\"a} }1
  {ö}{ {\"o} }1
  {~}{ {\textasciitilde} }1
  {а}{ {\selectfont\char224} }1
  {б}{ {\selectfont\char225} }1
  {в}{ {\selectfont\char226} }1
  {г}{ {\selectfont\char227} }1
  {д}{ {\selectfont\char228} }1
  {е}{ {\selectfont\char229} }1
  {ё}{ {\"e} }1
  {ж}{ {\selectfont\char230} }1
  {з}{ {\selectfont\char231} }1
  {и}{ {\selectfont\char232} }1
  {й}{ {\selectfont\char233} }1
  {к}{ {\selectfont\char234} }1
  {л}{ {\selectfont\char235} }1
  {м}{ {\selectfont\char236} }1
  {н}{ {\selectfont\char237} }1
  {о}{ {\selectfont\char238} }1
  {п}{ {\selectfont\char239} }1
  {р}{ {\selectfont\char240} }1
  {с}{ {\selectfont\char241} }1
  {т}{ {\selectfont\char242} }1
  {у}{ {\selectfont\char243} }1
  {ф}{ {\selectfont\char244} }1
  {х}{ {\selectfont\char245} }1
  {ц}{ {\selectfont\char246} }1
  {ч}{ {\selectfont\char247} }1
  {ш}{ {\selectfont\char248} }1
  {щ}{ {\selectfont\char249} }1
  {ъ}{ {\selectfont\char250} }1
  {ы}{ {\selectfont\char251} }1
  {ь}{ {\selectfont\char252} }1
  {э}{ {\selectfont\char253} }1
  {ю}{ {\selectfont\char254} }1
  {я}{ {\selectfont\char255} }1
  {А}{ {\selectfont\char192} }1
  {Б}{ {\selectfont\char193} }1
  {В}{ {\selectfont\char194} }1
  {Г}{ {\selectfont\char195} }1
  {Д}{ {\selectfont\char196} }1
  {Е}{ {\selectfont\char197} }1
  {Ё}{ {\"E} }1
  {Ж}{ {\selectfont\char198} }1
  {З}{ {\selectfont\char199} }1
  {И}{ {\selectfont\char200} }1
  {Й}{ {\selectfont\char201} }1
  {К}{ {\selectfont\char202} }1
  {Л}{ {\selectfont\char203} }1
  {М}{ {\selectfont\char204} }1
  {Н}{ {\selectfont\char205} }1
  {О}{ {\selectfont\char206} }1
  {П}{ {\selectfont\char207} }1
  {Р}{ {\selectfont\char208} }1
  {С}{ {\selectfont\char209} }1
  {Т}{ {\selectfont\char210} }1
  {У}{ {\selectfont\char211} }1
  {Ф}{ {\selectfont\char212} }1
  {Х}{ {\selectfont\char213} }1
  {Ц}{ {\selectfont\char214} }1
  {Ч}{ {\selectfont\char215} }1
  {Ш}{ {\selectfont\char216} }1
  {Щ}{ {\selectfont\char217} }1
  {Ъ}{ {\selectfont\char218} }1
  {Ы}{ {\selectfont\char219} }1
  {Ь}{ {\selectfont\char220} }1
  {Э}{ {\selectfont\char221} }1
  {Ю}{ {\selectfont\char222} }1
  {Я}{ {\selectfont\char223} }1
  {і}{ {\selectfont\char105} }1
  {ї}{ {\selectfont\char168} }1
  {є}{ {\selectfont\char185} }1
  {ґ}{ {\selectfont\char160} }1
  {І}{ {\selectfont\char73} }1
  {Ї}{ {\selectfont\char136} }1
  {Є}{ {\selectfont\char153} }1
  {Ґ}{ {\selectfont\char128} }1
  % \definecolor{brackets}{HTML}{B40000} % цвет скобок в коде
  % {\{}{ { {\color{brackets}\{} } }1 % Цвет скобок {
  % {\} }{ { {\color{brackets}\} } } }1 % Цвет скобок }
}

%--------------------------------------

% Table of contents
\usepackage{tocloft}
\usepackage{adjustbox}
\renewcommand{\cftpartleader}{\cftdotfill{\cftdotsep}} % for parts
\renewcommand{\cftsecleader}{\cftdotfill{\cftdotsep}} % for sections, if you really want! (It is default in report and book class (So you may not need it).

\usepackage{titlesec}
\titleformat{\paragraph} % Format the paragraph command
[hang] % Shape
{\normalfont\normalsize\bfseries} % Format for the text
{} % Label (we leave this empty to remove the number)
{0em} % Horizontal separation between label and title
{} % Code before the title body

% --- REDUCE PARAGRAPH INDENT IN TOC ---
% Reduce the space before the paragraph number (if numbered)
\setlength{\cftparaindent}{3em} % Default is often 2.5em or larger
% Reduce the box reserved for the paragraph number
% \setlength{\cftparanumwidth}{2.3em} % Default is often 3.0em
% \setlength{\cftparanumwidth}{0em}

% THIS IS THE CRUCIAL LINE: Set depth to 4 (paragraph) or even 5 (subparagraph)
\setcounter{tocdepth}{4} % 4 = paragraph, 5 = subparagraph
% --------------------------------------

%%% Математические пакеты %%%
%--------------------------------------
\usepackage{amsthm,amsfonts,amsmath,amssymb,amscd}  % Математические дополнения от AMS
\usepackage{mathtools}                              % Добавляет окружение multlined
\usepackage[perpage]{footmisc}
%--------------------------------------

%--------------------------------------
%			НАЧАЛО ДОКУМЕНТА
%--------------------------------------

\begin{document}

%--------------------------------------
%			ТИТУЛЬНЫЙ ЛИСТ
%--------------------------------------
\begin{titlepage}
\thispagestyle{empty}
\newpage


%Шапка титульного листа
%--------------------------------------
\vspace*{-60pt}
\hspace{-65pt}
\begin{minipage}{0.3\textwidth}
\hspace*{-20pt}\centering

\includegraphics[width=\textwidth]{emblem}
\end{minipage}
\begin{minipage}{0.67\textwidth}\small \textbf{
\vspace*{-0.7ex}
\hspace*{-6pt}\centerline{Министерство науки и высшего образования Российской Федерации}
\vspace*{-0.7ex}
\centerline{Федеральное государственное бюджетное образовательное учреждение }
\vspace*{-0.7ex}
\centerline{высшего образования}
\vspace*{-0.7ex}
\centerline{<<Московский государственный технический университет}
\vspace*{-0.7ex}
\centerline{имени Н.Э. Баумана}
\vspace*{-0.7ex}
\centerline{(национальный исследовательский университет)>>}
\vspace*{-0.7ex}
\centerline{(МГТУ им. Н.Э. Баумана)}}
\end{minipage}
%--------------------------------------

%Полосы
%--------------------------------------
\vspace{-25pt}
\hspace{-35pt}\rule{\textwidth}{2.3pt}

\vspace*{-20.3pt}
\hspace{-35pt}\rule{\textwidth}{0.4pt}
%--------------------------------------

\vspace{1.5ex}
\hspace{-35pt} \noindent \small ФАКУЛЬТЕТ\hspace{80pt} <<Информатика и системы управления>>

\vspace*{-16pt}
\hspace{47pt}\rule{0.83\textwidth}{0.4pt}

\vspace{0.5ex}
\hspace{-35pt} \noindent \small КАФЕДРА\hspace{50pt} <<Теоретическая информатика и компьютерные технологии>>

\vspace*{-16pt}
\hspace{30pt}\rule{0.866\textwidth}{0.4pt}
  
\vspace{11em}

\begin{center}
\Large {\bf Лабораторная работа № 3} \\ 
\large {\bf по курсу <<Базы данных>>} \\

\end{center}\normalsize

\vspace{8em}


\begin{flushright}
 \begin{minipage}{0.5\linewidth}
      {Студент: Федуков А. А. \hspace*{15pt}\\ 
      Группа: ИУ9-52Б \\ 
      Преподаватель: Вишняков И.Э. \hspace*{15pt}}
  \end{minipage}
\end{flushright}
\bigskip
\vfill
\begin{center}
\textsl{Москва, 2025 г.}
\end{center}
\end{titlepage}
%--------------------------------------
%		КОНЕЦ ТИТУЛЬНОГО ЛИСТА
%--------------------------------------

\renewcommand{\ttdefault}{pcr}

\setlength{\tabcolsep}{3pt}
\newpage
\setcounter{page}{2}

{
  \hypersetup{linkcolor=black}
  \tableofcontents
}


\addcontentsline{toc}{section}{Постановка задачи}
\section*{Постановка задачи}\label{Sect::goal}
Целью данной лабораторной работы является преобразование модели «сущность-связь» в реляционную модель.

Для достижения поставленной цели необходимо выполнить следующие задачи:
\begin{itemize}
    \item Преобразовать модель «сущность-связь», созданную в лабораторной работе №1, в реляционную модель согласно процедуре преобразования.
    \item Обосновать выбор типов данных, ключей, а также правил обеспечения ограничений минимальной кардинальности.
\end{itemize}

\addcontentsline{toc}{section}{Практическая реализация}
\section*{Практическая реализация}\label{Sect::implementation}





\paragraph{Диаграмма модели}
На рисунке \ref{fig:er_diagram} представлена ER-диаграмма, отображающая сущности, их атрибуты и связи, разработанная для предметной области «Поликлиника».


\begin{figure}[h!]
\centering
\includegraphics[width=0.9\textwidth]{Hospital ER Diagram}
\caption{ER-диаграмма системы «Поликлиника»}
\label{fig:er_diagram}
\end{figure}

\clearpage

\paragraph{Обоснование кардинальных чисел связей}
\begin{itemize}

    % Doctor ---
    \item \textbf{Appointment <--> Doctor} Один-ко-многим
    \begin{itemize}
        \item \textbf{Appointment --> Doctor} 1:1. Один прием может у одного и только одного врача
        \item \textbf{Appointment <-- Doctor} 0:N. Один доктор может вести несколько приемов, а может не иметь их вообще
    \end{itemize}

    \item \textbf{Prescription <--> Doctor} Один-ко-многим
    \begin{itemize}
        \item \textbf{Prescription --> Doctor} 1:1. Один рецепт может быть выписан только одним врачом.
        \item \textbf{Prescription <-- Doctor} 0:N. Один доктор может выписать несколько рецептов, но вообще не обязан их выписывать
    \end{itemize}

    \item \textbf{Assay <--> Doctor} Один-ко-многим
    \begin{itemize}
        \item \textbf{Assay --> Doctor} 0:1. Один анализ может быть проведен только одним врачом, но пока его нет, врач не назначен
        \item \textbf{Assay <-- Doctor} 0:N. Один доктор может провести несколько анализов, но вообще не обязан их проводить
    \end{itemize}


    % Patient ---
    \item \textbf{Appointment <--> Patient} Один-ко-многим
    \begin{itemize}
        \item \textbf{Appointment --> Patient} 1:1. Один прием может быть у только одного пациента
        \item \textbf{Appointment <-- Patient} 0:N. Один пациент может записаться на несколько приемов, но вообще не обязан иметь эти записи
    \end{itemize}

    \item \textbf{Prescription <--> Patient} Один-ко-многим
    \begin{itemize}
        \item \textbf{Prescription --> Patient} 1:1. Один рецепт может быть выписан только одному пациенту.
        \item \textbf{Prescription <-- Patient} 0:N. Один пациент может получить несколько рецептов, но вообще не обязан их получать
    \end{itemize}

    \item \textbf{Assay <--> Patient} Один-ко-многим
    \begin{itemize}
        \item \textbf{Assay --> Patient} 1:1. Один анализ может быть проведен только одному пациенту.
        \item \textbf{Assay <-- Patient} 0:N. Один пациент может пройти несколько анализов, но вообще не обязан их проходить
    \end{itemize}
\end{itemize}

\end{document}